\documentclass{article}
\usepackage[T1]{fontenc}
\usepackage{lmodern}
\usepackage{graphicx}
\usepackage{amsmath,amssymb}
\usepackage[a4paper,margin=2cm]{geometry}
\usepackage[absolute,overlay]{textpos}
\setlength{\TPHorizModule}{1mm}
\setlength{\TPVertModule}{1mm}
\usepackage[indonesian]{babel}
\usepackage{hyperref}
\setlength{\emergencystretch}{2em}
% unit: Halaman 37

\begin{document}

% === Halaman 37 (python/native) ===

Implementasi DES

• DES sudah diimplementasikan baik sebagai perangkat lunak maupun 

perangkat keras.

• Dalam bentuk perangkat keras, DES diimplementasikan di dalam chip. 

Setiap detik chip ini dapat mengenkripsikan 16,8 juta blok (atau 1 

gigabit per detik).

• Implementasi DES ke dalam perangkat lunak dapat melakukan 

enkripsi 32.000 blok per detik (dijalankan pada komputer mainframe 

IBM 3090, yaitu komputer tercepat saat itu / tahun 1976).

37

\end{document}
